\subsection{Módulo de Administración e Inteligencia de Negocio (ADM)}
\noindent\textbf{Ficha Técnica: ADM-013 - Gestión de Roles y Matriz de Permisos (RBAC Industrial)}
\footnotesize\setlength{\tabcolsep}{4pt}\renewcommand{\arraystretch}{1.2}
\rowcolors{2}{tableBlue}{white}\begin{longtable}{|p{0.21\textwidth}|p{0.74\textwidth}|}\hline
\textbf{Número} & ADM-013 \\ \hline
\textbf{Usuario} & Administrador del Sistema \\ \hline
\textbf{Prioridad} & MUST \\ \hline
\textbf{Puntos Estimados de Esfuerzo} & 5 \\ \hline
\textbf{Descripción} & \parbox[t]{\linewidth}{\vspace{2pt}\begin{itemize}[leftmargin=1.5em, nosep, topsep=0pt]
\item \textbf{Como} Administrador del Sistema
\item \textbf{Quiero} definir roles de usuario con una matriz de permisos granular (Crear, Editar, Eliminar, Aprobar, Cerrar, Auditar)
\item \textbf{Para} garantizar la segregación de funciones requerida por ISO 55001 y evitar que un mismo usuario solicite y apruebe acciones críticas.
\end{itemize}\vspace{2pt}} \\ \hline
\textbf{Observaciones} & Esta funcionalidad es fundamental para la seguridad del sistema y el cumplimiento de la norma ISO 55001 (5.3), la cual exige que la autoridad y responsabilidad estén explícitamente definidas. En esta fase se implementarán los roles base (Administrador, Planificador, Técnico y Auditor) con una matriz de permisos granular que asegure la segregación de funciones (SoD) y el principio de menor privilegio. \\ \hline
\end{longtable}\normalsize
\noindent\textit{Criterios de Aceptación:}
\footnotesize\begin{longtable}{|p{0.28\textwidth}|p{0.32\textwidth}|p{0.34\textwidth}|}\hline
\textbf{Contexto} & \textbf{Acción} & \textbf{Resultado} \\ \hline
\parbox[t]{\linewidth}{Dado que el Administrador necesita crear un perfil para ``Auditores Externos''.} & \parbox[t]{\linewidth}{Cuando crea el rol ``Auditor ISO'' y le asigna permisos de Solo Lectura en todas las áreas funcionales del sistema (Mantenimiento, Inventario, Visualización de Activos) pero ningún permiso de edición o borrado.} & \parbox[t]{\linewidth}{Entonces el sistema guarda el rol y permite asignarlo a usuarios, garantizando que estos puedan consultar toda la información pero no realizar modificaciones operativas.} \\ \hline
\parbox[t]{\linewidth}{Dado que se está configurando el rol de ``Planificador''.} & \parbox[t]{\linewidth}{Cuando el Administrador despliega los permisos del módulo ``Órdenes de Trabajo''.} & \parbox[t]{\linewidth}{Entonces debe ver opciones separadas para: ``Crear'', ``Editar'', ``Eliminar'', ``Aprobar Planeación'' y ``Cerrar Técnicamente''.} \\ \hline
\parbox[t]{\linewidth}{Dado que existe el rol ``Super Admin'' (Root).} & \parbox[t]{\linewidth}{Cuando el Administrador intenta eliminarlo o quitarle todos los permisos.} & \parbox[t]{\linewidth}{Entonces el sistema bloquea la acción con un mensaje de error crítico: ``No se puede modificar la integridad del rol raíz''.} \\ \hline
\end{longtable}\normalsize
\vspace{1em}
\noindent\textbf{Ficha Técnica: ADM-014 - Gestión de Usuarios y Ciclo de Vida Seguro}
\footnotesize\setlength{\tabcolsep}{4pt}\renewcommand{\arraystretch}{1.2}
\rowcolors{2}{tableBlue}{white}\begin{longtable}{|p{0.21\textwidth}|p{0.74\textwidth}|}\hline
\textbf{Número} & ADM-014 \\ \hline
\textbf{Usuario} & Administrador del Sistema \\ \hline
\textbf{Prioridad} & MUST \\ \hline
\textbf{Puntos Estimados de Esfuerzo} & 3 \\ \hline
\textbf{Descripción} & \parbox[t]{\linewidth}{\vspace{2pt}\begin{itemize}[leftmargin=1.5em, nosep, topsep=0pt]
\item \textbf{Como} Administrador del Sistema
\item \textbf{Quiero} gestionar el ciclo de vida completo de los usuarios (Crear, Suspender, Restaurar, Resetear Contraseña) de forma segura y auditable
\item \textbf{Para} proteger el acceso al Gemelo Digital y garantizar la trazabilidad de las acciones realizadas por exempleados.
\end{itemize}\vspace{2pt}} \\ \hline
\textbf{Observaciones} & Garantizar un acceso seguro es el primer paso para la integridad del sistema. Siguiendo la norma ISO 55001 (7.2), se gestionará el ciclo de vida completo de los usuarios incluyendo políticas anti-fuerza bruta y recuperación segura de contraseñas. La política de eliminación lógica (soft-delete) es mandatoria para preservar el historial completo de la auditoría. \\ \hline
\end{longtable}\normalsize
\noindent\textit{Criterios de Aceptación:}
\footnotesize\begin{longtable}{|p{0.28\textwidth}|p{0.32\textwidth}|p{0.34\textwidth}|}\hline
\textbf{Contexto} & \textbf{Acción} & \textbf{Resultado} \\ \hline
\parbox[t]{\linewidth}{Dado que el usuario ``Juan Pérez'' ha dejado la empresa pero tiene 50 Órdenes de Trabajo cerradas en su historial.} & \parbox[t]{\linewidth}{Cuando el Administrador hace clic en ``Desactivar Usuario''.} & \parbox[t]{\linewidth}{Entonces el sistema cambia su estado a INACTIVO, revoca su token de sesión actual, pero mantiene su nombre visible en todos los registros históricos (Auditoría Trail).} \\ \hline
\parbox[t]{\linewidth}{Dado que un usuario intenta iniciar sesión con una contraseña incorrecta 3 veces consecutivas.} & \parbox[t]{\linewidth}{Cuando falla el cuarto intento.} & \parbox[t]{\linewidth}{Entonces el sistema bloquea la cuenta temporalmente por 15 minutos y envía una alerta de seguridad al correo registrado.} \\ \hline
\parbox[t]{\linewidth}{Dado que un usuario olvidó su contraseña.} & \parbox[t]{\linewidth}{Cuando solicita el restablecimiento desde la pantalla de Login.} & \parbox[t]{\linewidth}{Entonces recibe un enlace seguro (token de un solo uso) con validez de 1 hora para definir una nueva credencial que cumpla la política de complejidad (mínimo 12 caracteres).} \\ \hline
\end{longtable}\normalsize
\vspace{1em}
\noindent\textbf{Ficha Técnica: ADM-015 - Configuración de Integración ERP (Simulación)}
\footnotesize\setlength{\tabcolsep}{4pt}\renewcommand{\arraystretch}{1.2}
\rowcolors{2}{tableBlue}{white}\begin{longtable}{|p{0.21\textwidth}|p{0.74\textwidth}|}\hline
\textbf{Número} & ADM-015 \\ \hline
\textbf{Usuario} & Administrador del Sistema \\ \hline
\textbf{Prioridad} & COULD \\ \hline
\textbf{Puntos Estimados de Esfuerzo} & 8 \\ \hline
\textbf{Descripción} & \parbox[t]{\linewidth}{\vspace{2pt}\begin{itemize}[leftmargin=1.5em, nosep, topsep=0pt]
\item \textbf{Como} Administrador del Sistema
\item \textbf{Quiero} configurar la conexión API con una instancia simulada de ERP (Mock Server u Odoo)
\item \textbf{Para} demostrar la sincronización automática de datos de inventario y asegurar que las requisiciones técnicas disparadas por la condición del activo se conviertan correctamente en procesos de compra financieros en el ERP.
\end{itemize}\vspace{2pt}} \\ \hline
\textbf{Observaciones} & Para validar la arquitectura de interoperabilidad industrial bajo estándares MIMOSA (OSA-EAI), se implementará una simulación funcional usando un servidor de pruebas o una instancia de código abierto. Esto permitirá demostrar la capacidad del sistema para sincronizar datos de inventario con plataformas externas de forma segura y escalable. \\ \hline
\end{longtable}\normalsize
\noindent\textit{Criterios de Aceptación:}
\footnotesize\begin{longtable}{|p{0.28\textwidth}|p{0.32\textwidth}|p{0.34\textwidth}|}\hline
\textbf{Contexto} & \textbf{Acción} & \textbf{Resultado} \\ \hline
\parbox[t]{\linewidth}{Dado que el Administrador está en la pantalla de ``Integraciones'' y ha ingresado una API Key válida del ERP.} & \parbox[t]{\linewidth}{Cuando hace clic en el botón ``Probar Conexión''.} & \parbox[t]{\linewidth}{Entonces el sistema debe mostrar un mensaje de éxito: ``Conexión con ERP establecida''.} \\ \hline
\parbox[t]{\linewidth}{Dado que el Administrador está en la pantalla de ``Integraciones''.} & \parbox[t]{\linewidth}{Cuando intenta probar la conexión con una API Key incorrecta o expirada.} & \parbox[t]{\linewidth}{Entonces el sistema debe mostrar un mensaje de error claro: ``Fallo la autenticación. Verifique la API Key.''} \\ \hline
\end{longtable}\normalsize
\vspace{1em}
\noindent\textbf{Ficha Técnica: ADM-016 - Publicación de Anuncios del Sistema}
\footnotesize\setlength{\tabcolsep}{4pt}\renewcommand{\arraystretch}{1.2}
\rowcolors{2}{tableBlue}{white}\begin{longtable}{|p{0.21\textwidth}|p{0.74\textwidth}|}\hline
\textbf{Número} & ADM-016 \\ \hline
\textbf{Usuario} & Administrador del Sistema \\ \hline
\textbf{Prioridad} & COULD \\ \hline
\textbf{Puntos Estimados de Esfuerzo} & 2 \\ \hline
\textbf{Descripción} & \parbox[t]{\linewidth}{\vspace{2pt}\begin{itemize}[leftmargin=1.5em, nosep, topsep=0pt]
\item \textbf{Como} Administrador del Sistema
\item \textbf{Quiero} publicar un anuncio global o por rol
\item \textbf{Para} comunicar a los usuarios sobre ventanas de mantenimiento u otras noticias importantes de la plataforma.
\end{itemize}\vspace{2pt}} \\ \hline
\textbf{Observaciones} & Se implementa un sistema de comunicación interna para notificar a los usuarios sobre eventos operativos o ventanas de mantenimiento. Es una funcionalidad de mejora de la experiencia de usuario que facilita la gestión del cambio y la difusión de información técnica relevante. \\ \hline
\end{longtable}\normalsize
\noindent\textit{Criterios de Aceptación:}
\footnotesize\begin{longtable}{|p{0.28\textwidth}|p{0.32\textwidth}|p{0.34\textwidth}|}\hline
\textbf{Contexto} & \textbf{Acción} & \textbf{Resultado} \\ \hline
\parbox[t]{\linewidth}{Dado que el Administrador está en la pantalla de ``Publicar Anuncio''.} & \parbox[t]{\linewidth}{Cuando escribe un mensaje, selecciona el rol ``Técnico'' como destinatario y lo publica.} & \parbox[t]{\linewidth}{Entonces solo los usuarios con el rol ``Técnico'' deben ver un banner con ese anuncio la próxima vez que inicien sesión.} \\ \hline
\end{longtable}\normalsize
\vspace{1em}
\noindent\textbf{Ficha Técnica: ADM-017 - Dashboard Ejecutivo de KPIs}
\footnotesize\setlength{\tabcolsep}{4pt}\renewcommand{\arraystretch}{1.2}
\rowcolors{2}{tableBlue}{white}\begin{longtable}{|p{0.21\textwidth}|p{0.74\textwidth}|}\hline
\textbf{Número} & ADM-017 \\ \hline
\textbf{Usuario} & Gerente de Planta \\ \hline
\textbf{Prioridad} & SHOULD \\ \hline
\textbf{Puntos Estimados de Esfuerzo} & 8 \\ \hline
\textbf{Descripción} & \parbox[t]{\linewidth}{\vspace{2pt}\begin{itemize}[leftmargin=1.5em, nosep, topsep=0pt]
\item \textbf{Como} Gerente de Planta
\item \textbf{Quiero} visualizar un Dashboard Ejecutivo con KPIs clave (MTBF, Costos, Disponibilidad) que permita filtrar por fecha y área
\item \textbf{Para} identificar tendencias de fallos y ``activos problema'' (Bad Actors) para dirigir el presupuesto de mantenimiento donde más se necesita.
\end{itemize}\vspace{2pt}} \\ \hline
\textbf{Observaciones} & El dashboard ejecutivo sintetiza la información operativa para la toma de decisiones estratégicas de presupuesto y confiabilidad. Basado en estándares ISO 14224 y mejores prácticas de MRO (Wireman Vol 2), el sistema permitirá identificar tendencias y activos críticos mediante un procesamiento eficiente de datos históricos. \\ \hline
\end{longtable}\normalsize
\noindent\textit{Criterios de Aceptación:}
\footnotesize\begin{longtable}{|p{0.28\textwidth}|p{0.32\textwidth}|p{0.34\textwidth}|}\hline
\textbf{Contexto} & \textbf{Acción} & \textbf{Resultado} \\ \hline
\parbox[t]{\linewidth}{Dado que el Gerente ha iniciado sesión.} & \parbox[t]{\linewidth}{Cuando navega al módulo de ``Inteligencia de Mantenimiento''.} & \parbox[t]{\linewidth}{Entonces el sistema debe mostrar los indicadores MTBF, MTTR y Disponibilidad del mes en curso para toda la planta, utilizando gráficos de tendencia (líneas) y medidores (gauges).} \\ \hline
\parbox[t]{\linewidth}{Dado que el Gerente está viendo el Dashboard General.} & \parbox[t]{\linewidth}{Cuando selecciona el rango ``Último Trimestre'' y filtra por la zona ``Casa de Bombas''.} & \parbox[t]{\linewidth}{Entonces todos los gráficos deben recalcularse automáticamente para mostrar solo los datos de esa zona en ese periodo.} \\ \hline
\parbox[t]{\linewidth}{Dado que el Gerente visualiza el gráfico de ``Top 10 Activos con más Paradas'' (Pareto).} & \parbox[t]{\linewidth}{Cuando hace clic en la barra del activo más crítico (el número 1).} & \parbox[t]{\linewidth}{Entonces el sistema debe redirigirlo a la ``Hoja de Vida'' de ese activo específico para ver el detalle de sus últimas fallas.} \\ \hline
\end{longtable}\normalsize
\vspace{1em}
\noindent\textbf{Ficha Técnica: ADM-018 - Exportación de Reportes en PDF}
\footnotesize\setlength{\tabcolsep}{4pt}\renewcommand{\arraystretch}{1.2}
\rowcolors{2}{tableBlue}{white}\begin{longtable}{|p{0.21\textwidth}|p{0.74\textwidth}|}\hline
\textbf{Número} & ADM-018 \\ \hline
\textbf{Usuario} & Gerente de Planta \\ \hline
\textbf{Prioridad} & SHOULD \\ \hline
\textbf{Puntos Estimados de Esfuerzo} & 5 \\ \hline
\textbf{Descripción} & \parbox[t]{\linewidth}{\vspace{2pt}\begin{itemize}[leftmargin=1.5em, nosep, topsep=0pt]
\item \textbf{Como} Gerente de Planta
\item \textbf{Quiero} descargar un reporte en PDF con los indicadores clave del mantenimiento (MTBF, MTTR, Disponibilidad)
\item \textbf{Para} analizar el desempeño mensual y presentarlo en las reuniones de gestión sin depender de la conexión al sistema.
\end{itemize}\vspace{2pt}} \\ \hline
\textbf{Observaciones} & La capacidad de exportación en PDF responde a la necesidad de análisis offline y presentación de resultados en reuniones de gestión. Se priorizará un diseño profesional del documento que incluya tablas de KPIs y gráficos de tendencia, garantizando que la información estratégica sea accesible sin conexión. \\ \hline
\end{longtable}\normalsize
\noindent\textit{Criterios de Aceptación:}
\footnotesize\begin{longtable}{|p{0.28\textwidth}|p{0.32\textwidth}|p{0.34\textwidth}|}\hline
\textbf{Contexto} & \textbf{Acción} & \textbf{Resultado} \\ \hline
\parbox[t]{\linewidth}{Dado que el Gerente está en la pantalla de ``Inteligencia de Negocio''.} & \parbox[t]{\linewidth}{Cuando selecciona el mes a consultar y hace clic en ``Descargar PDF''.} & \parbox[t]{\linewidth}{Entonces el sistema debe generar y descargar un archivo .pdf que contenga los gráficos de los indicadores seleccionados y mostrar el mensaje: ``Reporte generado correctamente.''} \\ \hline
\parbox[t]{\linewidth}{Dado que el Gerente intenta descargar el reporte del mes seleccionado.} & \parbox[t]{\linewidth}{Cuando el sistema detecta que no existen registros suficientes para calcular los indicadores.} & \parbox[t]{\linewidth}{Entonces debe mostrar un mensaje de error: ``No hay datos disponibles para generar el reporte del mes seleccionado.''} \\ \hline
\end{longtable}\normalsize
\vspace{1em}
\noindent\textbf{Ficha Técnica: ADM-019 - Sincronización de Datos con Herramientas BI}
\footnotesize\setlength{\tabcolsep}{4pt}\renewcommand{\arraystretch}{1.2}
\rowcolors{2}{tableBlue}{white}\begin{longtable}{|p{0.21\textwidth}|p{0.74\textwidth}|}\hline
\textbf{Número} & ADM-019 \\ \hline
\textbf{Usuario} & Gerente de Planta \\ \hline
\textbf{Prioridad} & COULD \\ \hline
\textbf{Puntos Estimados de Esfuerzo} & 5 \\ \hline
\textbf{Descripción} & \parbox[t]{\linewidth}{\vspace{2pt}\begin{itemize}[leftmargin=1.5em, nosep, topsep=0pt]
\item \textbf{Como} Gerente de Planta
\item \textbf{Quiero} exportar o sincronizar los datos de los KPIs directamente con herramientas externas (como Power BI o Excel)
\item \textbf{Para} tener presentaciones gerenciales siempre actualizadas y realizar análisis profundos sin tener que copiar y pegar gráficos manualmente cada mes.
\end{itemize}\vspace{2pt}} \\ \hline
\textbf{Observaciones} & Esta integración facilita el análisis profundo de datos al permitir el consumo de KPIs desde herramientas externas de inteligencia de negocio. Mediante una interfaz programática segura, el sistema asegura que los datos estratégicos estén disponibles para su uso en ecosistemas corporativos más amplios. \\ \hline
\end{longtable}\normalsize
\noindent\textit{Criterios de Aceptación:}
\footnotesize\begin{longtable}{|p{0.28\textwidth}|p{0.32\textwidth}|p{0.34\textwidth}|}\hline
\textbf{Contexto} & \textbf{Acción} & \textbf{Resultado} \\ \hline
\parbox[t]{\linewidth}{Dado que el Gerente está viendo el Dashboard de KPIs.} & \parbox[t]{\linewidth}{Cuando selecciona la opción ``Exportar Dataset estructurado''.} & \parbox[t]{\linewidth}{Entonces el sistema debe generar un archivo con formato compatible para ser importado en herramientas de oficina.} \\ \hline
\parbox[t]{\linewidth}{Dado que el Gerente utiliza una herramienta de BI externa.} & \parbox[t]{\linewidth}{Cuando configura el conector con la interfaz de conexión programática del sistema y su credencial.} & \parbox[t]{\linewidth}{Entonces la herramienta externa debe ser capaz de importar la tabla de datos históricos de KPIs y permitir su actualización automática.} \\ \hline
\end{longtable}\normalsize
\vspace{1em}
\noindent\textbf{Ficha Técnica: ADM-022 - Sugerencias de Optimización (PMO)}
\footnotesize\setlength{\tabcolsep}{4pt}\renewcommand{\arraystretch}{1.2}
\rowcolors{2}{tableBlue}{white}\begin{longtable}{|p{0.21\textwidth}|p{0.74\textwidth}|}\hline
\textbf{Número} & ADM-022 \\ \hline
\textbf{Usuario} & Supervisor de Mantenimiento \\ \hline
\textbf{Prioridad} & WON'T \\ \hline
\textbf{Puntos Estimados de Esfuerzo} & 5 \\ \hline
\textbf{Descripción} & \parbox[t]{\linewidth}{\vspace{2pt}\begin{itemize}[leftmargin=1.5em, nosep, topsep=0pt]
\item \textbf{Como} Supervisor de Mantenimiento
\item \textbf{Quiero} ver sugerencias automáticas sobre la frecuencia de mantenimiento basadas en el historial de fallas del activo
\item \textbf{Para} identificar equipos que están siendo sobre-mantenidos y ajustar el plan para reducir costos innecesarios.
\end{itemize}\vspace{2pt}} \\ \hline
\textbf{Observaciones} & Basado en el análisis de confiabilidad histórica, el sistema sugerirá optimizaciones al plan de mantenimiento para evitar el sobre-mantenimiento. Aunque su implementación completa requiere una base de datos histórica robusta, se define como una estrategia clave para la eficiencia de costos a largo plazo. \\ \hline
\end{longtable}\normalsize
\noindent\textit{Criterios de Aceptación:}
\footnotesize\begin{longtable}{|p{0.28\textwidth}|p{0.32\textwidth}|p{0.34\textwidth}|}\hline
\textbf{Contexto} & \textbf{Acción} & \textbf{Resultado} \\ \hline
\parbox[t]{\linewidth}{Dado que un activo tiene mantenimientos preventivos mensuales y cero fallas registradas en los últimos 12 meses.} & \parbox[t]{\linewidth}{Cuando el Supervisor accede al módulo de ``Optimización PMO''.} & \parbox[t]{\linewidth}{Entonces el sistema debe mostrar una tarjeta de sugerencia: ``Activo [ID] con alta confiabilidad. Se sugiere aumentar frecuencia a Bimestral''.} \\ \hline
\parbox[t]{\linewidth}{Dado que el sistema muestra una sugerencia de cambio de frecuencia.} & \parbox[t]{\linewidth}{Cuando el Supervisor hace clic en ``Aplicar Cambio''.} & \parbox[t]{\linewidth}{Entonces el Plan de Mantenimiento Preventivo asociado debe actualizarse automáticamente con la nueva frecuencia y registrar quién aprobó el cambio.} \\ \hline
\end{longtable}\normalsize
\vspace{1em}
\noindent\textbf{Ficha Técnica: ADM-024 - Configuración General del Sistema}
\footnotesize\setlength{\tabcolsep}{4pt}\renewcommand{\arraystretch}{1.2}
\rowcolors{2}{tableBlue}{white}\begin{longtable}{|p{0.21\textwidth}|p{0.74\textwidth}|}\hline
\textbf{Número} & ADM-024 \\ \hline
\textbf{Usuario} & Administrador del Sistema \\ \hline
\textbf{Prioridad} & WON'T \\ \hline
\textbf{Puntos Estimados de Esfuerzo} & 2 \\ \hline
\textbf{Descripción} & \parbox[t]{\linewidth}{\vspace{2pt}\begin{itemize}[leftmargin=1.5em, nosep, topsep=0pt]
\item \textbf{Como} Administrador del Sistema
\item \textbf{Quiero} configurar los parámetros globales del sistema (ej. Logo de la empresa, Zona Horaria, Unidades de Medida)
\item \textbf{Para} personalizar la plataforma según la identidad y ubicación de la planta sin necesidad de modificar el código.
\end{itemize}\vspace{2pt}} \\ \hline
\textbf{Observaciones} & Permite la personalización de los parámetros operativos y de identidad del sistema. La gestión de estas configuraciones globales asegura que la plataforma se adapte al entorno regional y a los estándares visuales de la organización sin requerir cambios a nivel de código. \\ \hline
\end{longtable}\normalsize
\noindent\textit{Criterios de Aceptación:}
\footnotesize\begin{longtable}{|p{0.28\textwidth}|p{0.32\textwidth}|p{0.34\textwidth}|}\hline
\textbf{Contexto} & \textbf{Acción} & \textbf{Resultado} \\ \hline
\parbox[t]{\linewidth}{Dado que el Administrador se encuentra en la pantalla de ``Configuración General''.} & \parbox[t]{\linewidth}{Cuando carga un archivo de imagen (PNG/JPG) en el campo ``Logo Corporativo'' y guarda los cambios.} & \parbox[t]{\linewidth}{Entonces el logo en la barra de navegación superior se actualiza inmediatamente para reflejar la nueva imagen.} \\ \hline
\parbox[t]{\linewidth}{Dado que el Administrador está en la configuración regional y la hora actual del sistema está en UTC.} & \parbox[t]{\linewidth}{Cuando selecciona ``UTC-5 (Colombia)'' en la lista desplegable de Zona Horaria.} & \parbox[t]{\linewidth}{Entonces todas las fechas y horas mostradas en las Órdenes de Trabajo deben ajustarse automáticamente a la nueva zona seleccionada.} \\ \hline
\end{longtable}\normalsize
\vspace{1em}
\noindent\textbf{Ficha Técnica: ADM-032 - Sistema de Auditoríaoría Inmutable (Auditoría Trail)}
\footnotesize\setlength{\tabcolsep}{4pt}\renewcommand{\arraystretch}{1.2}
\rowcolors{2}{tableBlue}{white}\begin{longtable}{|p{0.21\textwidth}|p{0.74\textwidth}|}\hline
\textbf{Número} & ADM-032 \\ \hline
\textbf{Usuario} & Gerente de Planta \\ \hline
\textbf{Prioridad} & MUST \\ \hline
\textbf{Puntos Estimados de Esfuerzo} & 5 \\ \hline
\textbf{Descripción} & \parbox[t]{\linewidth}{\vspace{2pt}\begin{itemize}[leftmargin=1.5em, nosep, topsep=0pt]
\item \textbf{Como} Gerente de Planta
\item \textbf{Quiero} consultar un registro cronológico e inalterable de todas las acciones críticas realizadas en el sistema (Creación, Modificación, Eliminación)
\item \textbf{Para} reconstruir la secuencia de eventos ante un incidente y demostrar cumplimiento normativo (ISO 55001 / ISO 14224).
\end{itemize}\vspace{2pt}} \\ \hline
\textbf{Observaciones} & La integridad de la información es un pilar de la norma ISO 55001. El sistema implementará un registro de auditoría inalterable que permita reconstruir cualquier incidente. Se utilizarán técnicas de encadenamiento de datos para asegurar que cada acción crítica sea trazable, cronológica y protegida contra manipulaciones. \\ \hline
\end{longtable}\normalsize
\noindent\textit{Criterios de Aceptación:}
\footnotesize\begin{longtable}{|p{0.28\textwidth}|p{0.32\textwidth}|p{0.34\textwidth}|}\hline
\textbf{Contexto} & \textbf{Acción} & \textbf{Resultado} \\ \hline
\parbox[t]{\linewidth}{Dado que un Planificador cambia la ``Prioridad'' de una Orden de Trabajo de ``Alta'' a ``Baja''.} & \parbox[t]{\linewidth}{Cuando guarda los cambios.} & \parbox[t]{\linewidth}{Entonces el sistema registra automáticamente una entrada de auditoría con: Quién: ``Carlos.Gomez''; Cuándo: 2026-02-21 10:30:05 UTC; Qué: Entidad ``WorkOrder'', ID ``WO-1052''; Detalle: Campo ``Priority'' cambió de ``Alta'' a ``Baja''.} \\ \hline
\parbox[t]{\linewidth}{Dado que el Auditoríaor está investigando el historial del activo ``Bomba P-101''.} & \parbox[t]{\linewidth}{Cuando filtra el registro de auditoría por ``EntityID = P-101''.} & \parbox[t]{\linewidth}{Entonces visualiza una línea de tiempo completa de todos los cambios de estado, mantenimientos y modificaciones de atributos desde su creación.} \\ \hline
\parbox[t]{\linewidth}{Dado que un usuario intenta eliminar un registro de ``Repuesto'' del inventario maestro.} & \parbox[t]{\linewidth}{Cuando confirma la acción.} & \parbox[t]{\linewidth}{Entonces el sistema marca el registro como inactivo (Soft-Delete) y crea una entrada de auditoría de tipo ``DELETE\_ATTEMPT'' con la justificación ingresada (si aplica).} \\ \hline
\end{longtable}\normalsize
\vspace{1em}